% ARCHIVO GENERADO por scripts/generar_capitulos.py. No editar a mano.
\section{Cambios de la iteración 2}




\begin{itemize}
  \item \textbf{Módulo nuevo \texttt{bitacora}} (DM-08, aprobado con el modelo v1.3).
  \item \textbf{Corrección de DM-08 detectada al dibujar el diagrama.} El texto aprobado dice que \texttt{bitacora} "depende de \texttt{core} y \texttt{cuentas}", pero \texttt{cuentas} también registra operaciones en la bitácora (operadores, horarios y autorizaciones), y eso formaría un \textbf{ciclo}. Para evitarlo, propongo que \texttt{bitacora} \textbf{solo dependa de \texttt{core}} y referencie al usuario con \texttt{settings.AUTH\_USER\_MODEL}, el mecanismo estándar de Django, sin importar \texttt{cuentas}. Así \texttt{cuentas} y \texttt{usuarios\_celdas} dependen de \texttt{bitacora}, y no al revés. Las tablas y las relaciones del modelo no cambian.
  \item \texttt{usuarios\_celdas} y \texttt{cuentas} dependen de \texttt{bitacora}; \texttt{usuarios\_celdas} también de \texttt{cuentas} (usuario que solicita o resuelve).
  \item ARQ-03 se actualizó con el módulo \texttt{bitacora} y las nuevas responsabilidades de \texttt{cuentas} y \texttt{usuarios\_celdas}.
\end{itemize}



\subsection{CHG-003: gestión de celdas (paquetes y secuencias)}
Solicitud de cambio aprobada el 2026-11-22, a partir de la observación del profesor.

\paragraph{Paquetes y secuencias}

\begin{itemize}
  \item \textbf{Sin módulos nuevos.} La gestión de celdas va en \texttt{usuarios\_celdas}. Para saber si una celda está ocupada o tiene historial usa la relación inversa \texttt{estadias}, igual que el mapa, sin importar \texttt{entradas\_salidas}. Así no se forma un ciclo.
  \item \textbf{Validación de ocupación:} el servicio consulta si la celda tiene una estadía abierta antes de inactivarla, cambiarle el tipo o eliminarla. Además, la clave foránea de la estadía impide en la base de datos eliminar una celda con historial.
  \item \textbf{Validación de código:} la aplicación valida el formato y la unicidad para dar un mensaje claro. La base de datos lo repite con un CHECK y la restricción única.
  \item \textbf{Secuencias:} SD-12 (crear celda) y SD-13 (inactivar celda) se agregaron al proyecto StarUML de la entrega.
\end{itemize}


